% Gerado por regressao/06_tabelas_latex.R (projeto Modelagem). Não editar à mão.
\begin{sidewaystable}
\caption{Resultados do modelo A (ciclo eleitoral) – efeitos fixos de município}\label{tab:regressao_principal_A}
\centering
\footnotesize
\setlength{\tabcolsep}{3pt}
\renewcommand{\arraystretch}{0.9}
\begin{tabular}{@{}>{\raggedright\arraybackslash}p{\dimexpr\linewidth-606pt\relax}>{\centering\arraybackslash}p{62pt}>{\centering\arraybackslash}p{60pt}>{\centering\arraybackslash}p{68pt}>{\centering\arraybackslash}p{62pt}>{\centering\arraybackslash}p{60pt}>{\centering\arraybackslash}p{60pt}>{\centering\arraybackslash}p{60pt}>{\centering\arraybackslash}p{60pt}>{\centering\arraybackslash}p{60pt}@{}}
\hline
\textbf{Variável} & \textbf{\hspace{0pt}Despesa total} & \textbf{\hspace{0pt}Saúde e Saneamento} & \textbf{\hspace{0pt}Educação e Cultura} & \textbf{\hspace{0pt}Habitação e Urbanismo} & \textbf{\hspace{0pt}Assistência Social} & \textbf{\hspace{0pt}Transporte} & \textbf{\hspace{0pt}Administração} & \textbf{\hspace{0pt}Agricultura} & \textbf{\hspace{0pt}Comunicações} \\ \hline
Ano eleitoral & 277,865\textsuperscript{***} & 53,163\textsuperscript{***} & 49,079\textsuperscript{*} & 126,814\textsuperscript{***} & 14,078\textsuperscript{***} & 41,741\textsuperscript{***} & 6,817 & 0,486 & $-$0,801\textsuperscript{***} \\
 & (102,741) & (10,703) & (29,440) & (31,902) & (4,274) & (10,585) & (10,968) & (2,026) & (0,218) \\
Ano pré-eleitoral & 162,806\textsuperscript{*} & 1,221 & 78,512\textsuperscript{**} & 34,083\textsuperscript{**} & 9,468\textsuperscript{**} & 8,342 & 10,347 & 5,168 & 0,656\textsuperscript{***} \\
 & (94,991) & (12,185) & (38,604) & (16,477) & (4,677) & (7,694) & (8,395) & (3,632) & (0,148) \\
Pandemia (2020) & $-$76,603 & 106,862\textsuperscript{***} & $-$242,843\textsuperscript{***} & 31,846\textsuperscript{*} & $-$0,062 & $-$0,516 & 22,273\textsuperscript{***} & 2,517 & 1,594\textsuperscript{***} \\
 & (66,107) & (7,028) & (21,771) & (19,287) & (2,979) & (8,917) & (6,793) & (2,523) & (0,295) \\
Prefeito na coalizão & 30,367\textsuperscript{***} & 4,917 & $-$0,971 & 3,187 & 0,665 & 2,588 & 6,870\textsuperscript{**} & $-$0,484 & $-$0,102 \\
do governador & (10,130) & (6,221) & (7,512) & (4,237) & (0,777) & (2,951) & (2,759) & (1,741) & (0,138) \\
Prefeito na coalizão & 67,005\textsuperscript{***} & 19,932\textsuperscript{***} & 31,782\textsuperscript{**} & 9,696\textsuperscript{*} & 4,793\textsuperscript{***} & $-$5,477 & 5,233\textsuperscript{*} & $-$1,503 & 0,442 \\
do presidente & (25,758) & (4,522) & (15,268) & (5,647) & (1,367) & (4,135) & (2,757) & (1,713) & (0,360) \\
Receita tributária & 0,934\textsuperscript{***} & 0,164\textsuperscript{***} & 0,208\textsuperscript{***} & 0,157\textsuperscript{***} & 0,030\textsuperscript{***} & 0,060\textsuperscript{***} & 0,169\textsuperscript{***} & 0,015\textsuperscript{***} & 0,002\textsuperscript{***} \\
per capita (R\$) & (0,089) & (0,010) & (0,025) & (0,032) & (0,003) & (0,014) & (0,018) & (0,003) & (0,001) \\
Transferências correntes & 0,638\textsuperscript{***} & 0,137\textsuperscript{***} & 0,163\textsuperscript{***} & 0,090\textsuperscript{***} & 0,023\textsuperscript{***} & 0,058\textsuperscript{***} & 0,095\textsuperscript{***} & 0,023\textsuperscript{***} & 0,002\textsuperscript{***} \\
per capita (R\$) & (0,057) & (0,010) & (0,026) & (0,007) & (0,002) & (0,008) & (0,007) & (0,003) & (0,000) \\
Proporção de jovens & $-$619,490 & 1.138,917\textsuperscript{***} & $-$3.453,893 & $-$1.527,928\textsuperscript{***} & 83,679 & 953,154\textsuperscript{***} & $-$851,194\textsuperscript{**} & $-$14,409 & $-$61,509\textsuperscript{***} \\
 & (2.842,870) & (307,711) & (2.291,827) & (531,047) & (159,415) & (330,930) & (352,276) & (140,987) & (19,745) \\
Proporção de idosos & 1.067,228 & 2.593,134\textsuperscript{**} & $-$13.602,563\textsuperscript{***} & 442,439 & $-$88,027 & 8.062,138\textsuperscript{***} & $-$969,992\textsuperscript{**} & 3.279,222\textsuperscript{***} & $-$92,944\textsuperscript{***} \\
 & (4.514,057) & (1.322,950) & (2.806,845) & (604,716) & (239,724) & (1.941,140) & (490,227) & (863,776) & (32,038) \\
Grau de urbanização & 1.095,956\textsuperscript{***} & $-$16,778 & 342,841\textsuperscript{***} & 269,831\textsuperscript{***} & 40,088\textsuperscript{***} & 493,564\textsuperscript{***} & $-$6,057 & 209,854\textsuperscript{***} & $-$14,643\textsuperscript{***} \\
 & (294,717) & (59,340) & (113,334) & (44,275) & (6,342) & (107,391) & (67,717) & (55,199) & (4,599) \\
População (log) & $-$1.562,263\textsuperscript{***} & $-$537,125\textsuperscript{***} & $-$737,046\textsuperscript{***} & 125,027 & $-$106,047\textsuperscript{***} & 181,138\textsuperscript{***} & $-$273,679\textsuperscript{***} & 44,459\textsuperscript{*} & $-$10,608\textsuperscript{**} \\
 & (461,765) & (84,603) & (183,921) & (77,769) & (14,577) & (62,684) & (40,086) & (24,767) & (4,139) \\
PIB nacional & 164,713 & 119,481\textsuperscript{***} & $-$126,888 & 150,102\textsuperscript{***} & 13,539 & 32,599 & 14,130 & 9,966 & $-$0,180 \\
(R\$ trilhões) & (267,297) & (23,512) & (124,729) & (51,626) & (11,896) & (20,734) & (24,579) & (12,426) & (0,444) \\
Tendência & $-$79,183 & 27,868\textsuperscript{**} & $-$90,719 & 30,180 & 2,390 & $-$29,110\textsuperscript{**} & $-$15,381 & $-$15,268\textsuperscript{*} & $-$0,964\textsuperscript{***} \\
 & (138,202) & (12,547) & (64,595) & (22,081) & (6,371) & (13,768) & (13,228) & (8,479) & (0,214) \\
Tendência ao quadrado & 9,260 & $-$1,154 & 12,117\textsuperscript{*} & $-$4,896\textsuperscript{**} & $-$0,161 & 0,148 & 0,512 & 0,347 & 0,034 \\
 & (13,826) & (1,292) & (6,203) & (2,441) & (0,624) & (1,125) & (1,319) & (0,670) & (0,022) \\
\hline
Observações & 70.669 & 70.526 & 70.550 & 69.574 & 70.485 & 54.423 & 70.553 & 64.464 & 12.360 \\
Municípios & 5.568 & 5.568 & 5.568 & 5.564 & 5.568 & 5.205 & 5.568 & 5.436 & 1.896 \\
R\textsuperscript{2} & 0,836 & 0,753 & 0,708 & 0,299 & 0,371 & 0,191 & 0,276 & 0,177 & 0,031 \\
\hline
\end{tabular}
\fonte{Elaboração própria.}
\nota{Variáveis dependentes: despesa per capita por grupo de funções (R\$ de 2025). Efeito fixo de município. Erros-padrão de Driscoll-Kraay entre parênteses. \textsuperscript{***}p$<$0,01; \textsuperscript{**}p$<$0,05; \textsuperscript{*}p$<$0,1. Os coeficientes de coalizão são apresentados para completude; a análise das coalizões usa o modelo B.}
\end{sidewaystable}
